\documentclass[11pt,a4paper]{article}
\usepackage{fontspec,amsmath,amssymb,amsthm,microtype}
\setmainfont{lmroman10-regular.otf}[BoldFont=lmroman10-bold.otf,ItalicFont=lmroman10-italic.otf,BoldItalicFont=lmroman10-bolditalic.otf]
\setmonofont{lmmono10-regular.otf}[Scale=MatchLowercase]
\usepackage[margin=25mm]{geometry}
\usepackage[unicode,hidelinks]{hyperref}
\setlength{\emergencystretch}{3em}
\setlength{\parskip}{0.4em}
\setlength{\parindent}{0pt}
\allowdisplaybreaks[2]
\raggedbottom
\title{\texorpdfstring{The supplied KIPLMC2 rate with an explicit dimension-\allowbreak{}independent constant}{The supplied KIPLMC2 rate with an explicit dimension-independent constant}}
\author{LangevinLean mathematical exposition}
\date{}
\begin{document}\maketitle
\noindent Source module: \texttt{KIPLMC/\allowbreak{}OBABOPaper\allowbreak{}Rate.\allowbreak{}lean}. The accompanying HTML pairs this mathematical exposition with the exact Lean statements and records the scope of the formalization.


The preceding actual sampler theorem gives a stronger intensive bias. This module derives the precise coarser form printed in the supplied manuscript, including its additional transient term. The implicit constant is replaced by an explicit expression that depends only on \(m,L,\gamma\), and the final result concerns the actual parameter-error integral.


\section{\texorpdfstring{Turn the square-\allowbreak{}root geometric factor into the printed exponential}{Turn the square-root geometric factor into the printed exponential}}


Let \(r=\sqrt{1-mh/(3\gamma)}\). The scalar bound \(1-x\le e^{-x}\), followed by taking square roots, gives \(r\le e^{-mh/(6\gamma)}\). Raising to the actual integer iteration index proves \(r^n\le e^{-mhn/(6\gamma)}\). The proof also handles a nonpositive square-root argument directly, so the scalar envelope itself does not need an extra admissibility hypothesis.



{\raggedright\noindent Lean declarations: \texttt{kiplmc2\_\allowbreak{}rate\_\allowbreak{}exp\_\allowbreak{}bound}.\par}

\section{\texorpdfstring{Make the hidden constant explicit}{Make the hidden constant explicit}}


The allowed step satisfies \(h\le m/(33\gamma^3)\). Consequently the physical bias coefficient is bounded by

\[B=\frac{6\sqrt2\gamma}{m}\left(4\sqrt L+\frac{Lm}{66\gamma^3}\right).\]

Choose \(C=\sqrt2+B+1/\sqrt m\). All three summands are nonnegative, so \(C\) dominates the transient coefficient, bias coefficient and equilibrium concentration coefficient separately. Lean proves these comparisons from positivity, without permitting dependence on the particle count, dimensions or step.



{\raggedright\noindent Lean declarations: \texttt{kiplmc2\allowbreak{}Paper\allowbreak{}Bias\allowbreak{}Coefficient}, \texttt{kiplmc2\allowbreak{}Paper\allowbreak{}Constant}, \texttt{kiplmc2\allowbreak{}Paper\allowbreak{}Constant\_\allowbreak{}bounds}.\par}

\section{\texorpdfstring{Compare each term with the supplied coarse particle envelope}{Compare each term with the supplied coarse particle envelope}}


For \(N\ge1\), \(D\le\sqrt N D+hN\sqrt{d_z}\) and \(\sqrt{d_z/N}\le\sqrt{Nd_z}\). The equilibrium term satisfies \(\sqrt{d_\theta/(Nm)}=(1/\sqrt m)\sqrt{d_\theta/N}\). Apply the exponential rate envelope, dominate each coefficient by \(C\), and add the three inequalities. This proves the exact printed shape, including the harmless nonnegative extra transient contribution \(hN\sqrt{d_z}\), from the sharper estimate.



{\raggedright\noindent Lean declarations: \texttt{kiplmc2\_\allowbreak{}paper\_\allowbreak{}envelope}.\par}

\section{\texorpdfstring{Apply the envelope to the actual particle parameter law}{Apply the envelope to the actual particle parameter law}}


The final theorem starts from the unconditional actual KIPLMC2 RMS estimate and supplies every scalar hypothesis using the original curvature, friction and step restrictions. With \(D=W_2(\nu,\pi_N)\), it proves

\[\sqrt{\mathbb E\lVert\theta_n-\theta_\star\rVert^2}\le C\left[(\sqrt ND+hN\sqrt{d_z})e^{-mhn/(6\gamma)}+h\sqrt{Nd_z}+\sqrt{d_\theta/N}\right].\]

This is the supplied manuscript theorem with an explicit constant and the actual independent Gaussian algorithm law. The proof does not assume its desired RMS estimate as an input.



{\raggedright\noindent Lean declarations: \texttt{kiplmc2\_\allowbreak{}parameter\_\allowbreak{}paper\_\allowbreak{}bound}.\par}

\end{document}
